
\section{Differentialgleichungen mit trennbaren Variablen (+)}\label{MA2-trennbare-differentialgleichungen}

\begin{table}[H]
  \centering
  \begin{tabularx}{\textwidth}{@{}>{\raggedright\arraybackslash}X r@{}}
    \toprule
    \multicolumn{2}{c}{\textbf{�bersicht �ber Lernziele und Voraussetzungen}} \\
    \midrule
    \textbf{Lernziele} & \\[.5\normalbaselineskip]
    \tabitem Differentialgleichungen mit trennbaren Variablen durch Separation der Variablen l�sen & \\
    \tabitem Anfangswertprobleme dieses Typs vollst�ndig bestimmen & \\[.5\normalbaselineskip]
    \midrule
    \textbf{Voraussetzungen} &  \\[.5\normalbaselineskip]
    \tabitem Integrale sicher berechnen k�nnen & \abschnittsverweis{MA2-integrationsregeln}  \\[.5\normalbaselineskip]
    \bottomrule
  \end{tabularx}
\end{table}

\noindent \ifintegriereaufgaben Die �bungsaufgaben zum Thema finden Sie \hyperref[MA2-trennbare-differentialgleichungen-aufgaben]{hier}.
\vspace{0.3cm}
\hrule
\fi


\subsection{Trennbare Variablen}

Die sogenannte "\,Trennung der Variablen"\, [separation of variables] ist ein beliebter Trick, um Anfangswertprobleme erster Ordnung (und nur erster Ordnung) zu l�sen. Differentialgleichungen erster Ordnung haben trennbare Variablen [separable variables], wenn gilt:
\begin{defbox}{Differentialgleichung mit trennbaren Variablen}{ }
	Die Differentialgleichung l��t sich eine Form bringen, in der
	\begin{itemize}
		\item auf der linken Seite das Produkt der Ableitung $y'$ mit einer Funktion $h$ steht, wobei $h$ nur von $y$ abh�ngt,
		\item auf der rechten Seite nur eine Funktion $f$ der unabh�ngigen Variable $x$ auftaucht.
	\end{itemize}
	Schematisch sieht sie dann so aus:
	\[ h(y(x)) \cdot y'(x) = f(x)\,.\]
\end{defbox}

\subsection{L�sungsstrategie}

Die grunds�tzliche Idee ist beide Seiten der Gleichung zu integrieren und die \underline{Substitutionsregel}\randverweis{MA2-integrationsregeln} anzuwenden. Die Substitution lautet
\[
	s = y(x) \Rightarrow \mathrm{d}s = y'(x)\, \mathrm d x\,.
\]
Anstatt $s$ schreibt man einfach $y$, also
\[
	y = y(x) \Rightarrow \mathrm{d}y = y'(x)\, \mathrm d x\,. \quad(*)
\]

Wir schauen uns das Vorgehen f�r das folgende Anfangswertproblem an:
\[
\begin{cases}
	y(3) =  7\,, \quad & x=3\,,\\[1ex]
	\displaystyle{e^{y(x)}} y'(x) = x^5, \quad & x > 3\,.
\end{cases}
\]
Wir integrieren beide Seiten der Differentialgleichung vom Startpunkt $(x,y)=(3,7)$ zu einem beliebigen Endpunkt $(x_1,y_1)$ und integrieren durch Substitution:
\begin{eqnarray*}
	\begin{array}{l rcl c rcl}
		& \displaystyle{\int\limits_{3}^{x_1}} e^y \displaystyle{\frac{\mathrm{d}y}{\mathrm{d}x}}\,\mathrm{d}x &=& \displaystyle{\int\limits_{3}^{x_1}} x^5 \, \mathrm{d}x  &\stackrel{\displaystyle{ (*)}}{\Leftrightarrow} & \displaystyle{\int\limits_{7}^{y_1}} e^y \,\mathrm{d}y &=& = \displaystyle{\int\limits_{3}^{x_1}} x^5 \, \mathrm{d}x \\[2ex]
		& & &  &\Leftrightarrow  & \Big[ \displaystyle{e^{y}}\Big]_{7}^{y_1} &=& \Big[ \displaystyle{\frac16 x^6}\Big]_{3}^{x_1} \\
		&&&&\Leftrightarrow  & \displaystyle{e^{y_1} - e^{7}} &=& \displaystyle{7\frac{x_1^6}{6} - \frac{3^6}{6}}
\end{array}
		\end{eqnarray*}
Und es ergibt sich die L�sung:
\[
 y_1 =\ln\left( e^7 + \frac{x_1^6}{6} - \frac{3^6}{6} \right) \,.
\]

\noindent Bemerkungen:
\begin{itemize}
	\item Das Aufl�sen nach $y_1$ im letzten Schritt kann problematisch sein - denken Sie an den Definitionsbereich der \underline{Umkehrfunktionen}\randverweis{MA1-eigenschaften-von-funktionen}.
	\item Der Name der L�sungmethode kommt daher, dass man die Differentialgleichung formal umformen kann zu
		\[
			e^y y' = x^5 \Leftrightarrow e^y \frac{\mathrm{d}y}{\mathrm{d}x} = x^5 \Rightarrow  e^y \mathrm{d}y = x^5 \mathrm{d}x\,.
	\]
In dieser Form sind die Variablen $x$ und $y$ getrennt und dann auf beiden Seiten von $(x,y)=(3,7)$ bis $(x_1,y_1)$ integriert:
			\[
			\int\limits_{7}^{y_1} e^y \,\mathrm{d}y = \int\limits_{3}^{x_1} x^5 \, \mathrm{d}x\,.
		\]
	\item Manche Leute verzichten auf die Anfangs- und Endwerte und schreiben unbestimmte Integrale, um die allgemeine L�sung der Differentialgleichung zu erhalten. Aber dann darf man die Integrationskonstante nicht vergessen:
			\[
			\int e^y \,\mathrm{d}y = \int x^5 \, \mathrm{d}x \Leftrightarrow e^y  = \frac{x^6}{6} + c\,.
		\]
	Die allgemeine L�sung lautet mit $C = \ln(c)$ also
		\[
			y(x)  = C\cdot \ln\left(\frac{x^6}{6}\right)\,.
	\]
\end{itemize}

\begin{exabox}{Nichtlinear von erster Ordnung}{ }
Anfangswertproblems:
	\[
		\begin{cases}
			y(0) = 1\,, &\quad x=0\,,\\
			y'(x)+y^2(x)  = 0\,, &\quad x>0\,.
		\end{cases}\]
Trennung der Variablen:
\begin{eqnarray*}
y'+y^2 = 0 \Leftrightarrow \frac{\mathrm{d}y}{\mathrm{d}x}  -y^2 &=&  0  \\
\stackrel{y\not=0}{\Leftrightarrow} \frac{1}{y^2}\cdot \frac{\mathrm{d}y}{\mathrm{d}x} &=& -1 \\
\Leftrightarrow \frac{1}{y^2}\cdot \mathrm{d}y &=& -\mathrm{d}x
\end{eqnarray*}
Integration:
\[
\int \frac{1}{y^2}\,\mathrm{d}y = -\int\,\mathrm{d}x  \stackrel{x\not=C}{\Leftrightarrow} y = \frac{1}{x-C} \,.
\]
Bestimmung der Integrationskonstanten:
\[ y(0) = -\frac{1}{C} \stackrel{!}{=} 1 \Leftrightarrow C = -1 \,.\]
L�sung:
\[ y(x) = \frac{1}{x+1}\,.\]
\end{exabox}

Bei den folgenden Differentialgleichungen f�hrt eine Substitution auf eine Form, die mit dem Schema dieses Abschnitts zu l�sen ist:

	\begin{center}
		\begin{tabular}{c|c|c}
			vorher & Substitution & nachher \\ \hline
			$y' =\tilde h(ax+by+c)$ & $ u = ax+by+c$ & $u' = a+b \tilde h(u)$ \\
			\hline
			$y' =\tilde h \left(\displaystyle{\frac{y}{x}}\right)$ & $ u = \displaystyle{\frac{y}{x}}$ & $u' = \displaystyle{\frac{\tilde h(u) - u}{x}}$ \\
		\end{tabular}
	\end{center}
